%\documentclass[12pt,a4paper]{article}  % Use this line if this document will be released
\documentclass[12pt,a4paper,draft]{beamer}  % Use this line if this document is a draft
\usepackage{ifdraft}


%% Bibliography
\usepackage{etoolbox}
\newcommand{\bibfile}{\jobname.bib}  % Name of the BibTeX file.
% ref.bib should be a symbolic link to the universal BibTeX file, which should be a local copy of
% https://github.com/equipez/bibliographie/blob/main/ref.bib
% Run `getbib` in the current directory under the draft mode to get the BibTeX file containing only
% the cited references. The name will be xyz.bib if this TeX file is xyz.tex.
\newcommand{\universalbib}{ref.bib}
\ifdraft{\IfFileExists{\universalbib}{\renewcommand{\bibfile}{\universalbib}}{}}{}
% The counter `cite' is used to count the number of citations.
\newcounter{cite}
\pretocmd{\cite}{\stepcounter{cite}}{}{}


%% Add line numbers in draft mode
\RequirePackage[mathlines]{lineno}
\ifdraft{\linenumbers}{}
\renewcommand{\linenumberfont}{\normalfont\scriptsize\sffamily\color{gray}}
\setlength{\linenumbersep}{\marginparsep}


%% Geometry
%\voffset=-1.5cm \hoffset=-1.4cm \textwidth=16cm \textheight=22.0cm  % Luis' setting
\usepackage[a4paper, textwidth=16.0cm, textheight=22.0cm]{geometry}
\renewcommand{\baselinestretch}{1.2}


%% Basic packages
\usepackage{amsmath,amsthm,amssymb,amsfonts}
\usepackage{mathtools}  % Provides \coloneqq
\usepackage{empheq}
\usepackage{xcolor}
\usepackage[bbgreekl]{mathbbol}
\DeclareSymbolFontAlphabet{\mathbbm}{bbold}
\DeclareSymbolFontAlphabet{\mathbb}{AMSb}
\usepackage{bbm}
\usepackage{upgreek}
\usepackage{accents}
\usepackage{xspace}
\usepackage{rotating}
\usepackage{multirow,booktabs}
\usepackage[en-US]{datetime2}


%% Format of the table of content
\usepackage[normalem]{ulem}
\usepackage[toc,page]{appendix}
\renewcommand{\appendixpagename}{\Large{Appendix}}
\renewcommand{\appendixname}{Appendix}
\renewcommand{\appendixtocname}{Appendix}
%\usepackage{sectsty}
\setcounter{tocdepth}{2}


%% Section title style
\usepackage{sectsty}
\sectionfont{\large}
\subsectionfont{\large}


%% Some colors
\definecolor{darkblue}{rgb}{0,0.1,0.5}
\definecolor{darkgreen}{rgb}{0,0.5,0.1}
\definecolor{darkyellow}{rgb}{0.65,0.65,0.01}


%% Todo notes
\ifdraft{
    \setlength{\marginparwidth}{2.42cm}
    \usepackage[tickmarkheight=3pt,textsize=small,backgroundcolor=blue!16,linecolor=purple,bordercolor=purple]{todonotes}
}{
    \newcommand{\todo}[1]{}
    \newcommand{\listoftodos}{}
}


%% Graph, tikz and pgf
%\usepackage{subfigure}
\setlength{\unitlength}{1mm}
% The \unitlength command is a Length command. It defines the units used in the Picture Environment.
\usepackage{graphicx}
%\usepackage{tikz,tikzscale,pgf,pgfarrows,pgfnodes,filecontents,tikz-cd}
\usepackage{tikz,tikzscale,pgf}
\usetikzlibrary{arrows,arrows.meta,patterns,positioning,decorations.markings,shapes}
\usepackage{pgfplots}
\usepackage{pgfplotstable}
\usepackage[justification=centering]{caption}
\usepgfplotslibrary{fillbetween}
\pgfplotsset{compat=1.11}


%% Turn off some unharmful warnings in draft mode
%% N.B.: DO NOT use `silence` together with `hyperref`. They will cause an infinite loop.
\ifdraft{
    \usepackage{silence}
    \WarningFilter{xcolor}{Incompatible color definition on}
    \WarningFilter{hyperref}{Draft mode on}
    \WarningFilter{refcheck}{Unused label}
    \WarningFilter{microtype}{`draft' option active}
    \WarningFilter{latex}{Writing or overwriting file} % Mute the warning about 'writing/overwriting file'
    \WarningFilter{latex}{Writing file} % Mute the warning about 'writing/overwriting file'
    \WarningFilter{latex}{Tab has} % Mute the warning about 'Tab has been converted to Blank Space'
    \WarningFilter{latex}{Marginpar on page} % Mute the warning about 'Marginpar on page xx moved'
    \WarningFilter{latex}{author given} % Mute the warning about 'No \author given'
}{}


%% Hyperref, url, and email
%% N.B.: DO NOT use `silence` together with `hyperref`. They will cause an infinite loop.
\ifdraft{\usepackage{refcheck}\newcommand{\url}{\texttt}}{
    \usepackage{hyperref}
    \hypersetup{colorlinks, linkcolor=darkblue, anchorcolor=darkblue, citecolor=darkblue, urlcolor=darkblue}
    \usepackage{url}
} % Check unused labels
\newcommand{\email}{\texttt}


%% Enumerate and itemize
\usepackage{eqlist}
\usepackage{enumitem}
\setlist[itemize]{leftmargin=*}
\setlist[enumerate]{leftmargin=*,label=\normalfont{(\alph*)}}


%% Algorithm environment
\usepackage[section]{algorithm}
\usepackage{algpseudocode,algorithmicx}
\newcommand{\INPUT}{\textbf{Input}}
\newcommand{\FOR}{\textbf{For}~}
\algrenewcommand\algorithmicrequire{\textbf{Input:}}
\algrenewcommand\algorithmicensure{\textbf{Output:}}
\algrenewcommand\alglinenumber[1]{\normalsize #1.}
\newcommand*\Let[2]{\State #1 $=$ #2}


%% Theorem-like environments
\newtheorem{theorem}{Theorem}[section]
\newtheorem{conjecture}{Conjecture}[section]
\newtheorem{corollary}{Corollary}[section]
\newtheorem{exercise}{Exercise}[section]
\newtheorem{lemma}{Lemma}[section]
\newtheorem{problem}{Problem}[section]
\newtheorem{proposition}{Proposition}[section]
\newtheorem{assumption}{Assumption}[section]
\newtheorem{example}{Example}[section]
\newtheorem{question}{Question}[section]
% Change theoremstyle to ``definition'', which uses textnormal for the text.
\theoremstyle{definition}
\newtheorem{definition}{Definition}[section]
\newtheorem{remark}{Remark}[section]
% proof
\usepackage{xpatch}
\xpatchcmd{\proof}{\itshape}{\normalfont\proofnamefont}{}{}
\newcommand{\proofnamefont}{\bfseries}

%% Equation numbering
\numberwithin{equation}{section}


%% Fine tuning
\usepackage{microtype}
\usepackage[nobottomtitles*]{titlesec} % No section title at the bottom of pages
% Prevent footnote from running to the next page
\interfootnotelinepenalty=10000
% No line break in inline math
\interdisplaylinepenalty=10000
\relpenalty=10000
\binoppenalty=10000
% No widow or orphan lines
\clubpenalty=10000
\widowpenalty=10000
\displaywidowpenalty=10000


% Use @ to put 1 math unit (mu) in math
% See https://nhigham.com/2013/01/07/fine-tuning-spacing-in-latex-equations/
% and also TeXbook p. 155.
\mathcode`@="8000{\catcode`\@=\active\gdef@{\mkern1mu}}


%% Operators, commands
\usepackage{relsize}
\usepackage{nccmath}
%\DeclareMathOperator*{\mcap}{\,\medmath{\bigcap}\,}
%\DeclareMathOperator*{\mcup}{\,\medmath{\bigcup}\,}
\DeclareMathOperator*{\mcap}{\,\mathsmaller{\bigcap}\,}
\DeclareMathOperator*{\mcup}{\,\mathsmaller{\bigcup}\,}
%\renewcommand{\cap}{\mcap}
%\renewcommand{\cup}{\mcup}

\newcommand{\ceil}[1]{ {\lceil{#1}\rceil} }
\newcommand{\floor}[1]{ {\lfloor{#1}\rfloor} }

\DeclareMathOperator{\tr}{tr}
\DeclareMathOperator{\sort}{sort}
\DeclareMathOperator*{\Argmax}{Argmax}
\DeclareMathOperator*{\Argmin}{Argmin}
\DeclareMathOperator*{\Arglocmin}{Arglocmin}
\DeclareMathOperator*{\argmax}{argmax}
\DeclareMathOperator*{\argmin}{argmin}
\DeclareMathOperator*{\diag}{diag}
\DeclareMathOperator*{\Diag}{Diag}
\DeclareMathOperator{\Span}{span}
\DeclareMathOperator{\med}{med}
\DeclareMathOperator{\essinf}{essinf}
\DeclareMathOperator{\cl}{cl}
\DeclareMathOperator{\vol}{vol}
\DeclareMathOperator{\comp}{C}
\DeclareMathOperator{\sign}{sign}
\DeclareMathOperator{\rank}{rank}
\DeclareMathOperator{\range}{range}
\DeclareMathOperator{\card}{card}
\DeclareMathOperator{\diam}{diam}
\DeclareMathOperator{\dist}{dist}
\newcommand{\disth}{{\operatorname{\updelta_{\sss{H}}}}}
\newcommand{\ind}{\mathbbm{1}}
%\newcommand*{\defeq}{\stackrel{\mbox{\normalfont\tiny{\textnormal{def}}}}{=}}
\newcommand\defeq{\mathrel{\overset{\makebox[0pt]{\mbox{\normalfont\tiny\sffamily def}}}{=}}}

\newcommand{\RR}{\mathbb{R}}
\newcommand{\BB}{\mathcal{B}}
\renewcommand{\SS}{\mathbb{S}}
\newcommand{\TT}{\mathcal{T}}
\newcommand{\ZZ}{\mathbb{Z}}
\newcommand{\NN}{\mathbb{N}}
\newcommand{\FF}{\mathcal{F}}
\newcommand{\CC}{\mathbb{C}}
\newcommand{\XX}{\mathcal{X}}
\newcommand{\sset}{\mathcal{S}}
\newcommand{\pen}{h}
\newcommand{\penpar}{\mu}
\newcommand{\res}{\rho}
\newcommand{\col}{r}
\newcommand{\ofd}{\mathcal{F}}
\newcommand{\stf}[1]{\mathbb{S}^{#1}}
\newcommand{\sss}[1]{{\scriptscriptstyle{#1}}}
\newcommand{\sK}{{\scriptscriptstyle{K}}}
\newcommand{\sT}{{\scriptscriptstyle{T}}}
\newcommand{\fro}{{\scriptstyle{\textnormal{F}}}}
\newcommand{\trs}{{\scriptstyle{\mathsf{T}}}}
\newcommand{\hmt}{{\scriptstyle{{\mathsf{H}}}}}
\newcommand{\pin}{{\scriptstyle{{\mathsf{+}}}}}
\newcommand{\inv}{{-1}}
\newcommand{\adj}{*}
\newcommand{\ones}{\mathbf{1}}

\newcommand{\cs}{\text{c}}
\newcommand{\hp}{\circ}
\newcommand{\cc}{\sss{\textnormal{C}}}
\newcommand{\dec}{\sss{\textnormal{D}}}
\newcommand{\cauchy}{\sss{\textnormal{C}}}
\newcommand{\scauchy}{\sss{\textnormal{S}}}
\newcommand{\crit}{\textnormal{crit}}
\newcommand{\rsg}{\hat{\partial}}
\newcommand{\gsg}{\partial}
\newcommand{\dom}{\textnormal{dom}}
\newcommand{\tf}{{\textnormal{f}}}
\newcommand{\tg}{{\textnormal{g}}}
\newcommand{\ts}{{\textnormal{s}}}
\newcommand{\st}{\textnormal{s.t.}}
\newcommand{\etc}{{etc.}\xspace}
\newcommand{\ie}{{i.e.}\xspace}
\newcommand{\eg}{{e.g.}\xspace}
\newcommand{\etal}{{et al.}\xspace}
\newcommand{\iid}{\text{i.i.d.}\xspace}
\newcommand{\as}{\text{a.s.}\xspace}

\newcommand{\me}{\mathrm{e}}
\newcommand{\md}{\mathrm{d}}
\newcommand{\mi}{\mathrm{i}}
\newcommand{\lev}{\mathrm{lev}}
\newcommand{\bA}{\mathbf{A}}
\newcommand{\bx}{\mathbf{u}}
%\newcommand{\bb}{\mathbf{f}}
\newcommand{\bb}{\mathbf{r}}
\newcommand{\nov}{n_{\textnormal{o}}}
\xspaceaddexceptions{]\}}
% tex.stackexchange.com/questions/15252/why-does-xspace-behave-differently-for-parenthesis-vs-braces-brackets
\newcommand{\MATLAB}{\textsc{Matlab}\xspace}
\newcommand{\octave}{\mbox{GNU Octave}\xspace}
\newcommand{\prblm}{\texttt}
\DeclareMathAlphabet{\mathsfit}{T1}{\sfdefault}{\mddefault}{\sldefault}
\SetMathAlphabet{\mathsfit}{bold}{T1}{\sfdefault}{\bfdefault}{\sldefault}
\newcommand{\prbb}{\mathsfit{p}}
\newcommand{\pp}{\mathsf{p}}
\newcommand{\qq}{\mathsf{q}}
\newcommand{\ttt}{\mathsfit{t}}
\newcommand{\tol}{\varepsilon}
\newcommand{\bt}{\mathbf{t}}
\newcommand{\br}{\mathbf{r}}
\newcommand{\dd}{\mathbf{d}}
\newcommand{\ii}{\mathbf{i}}
\newcommand{\jj}{\mathbf{j}}
\newcommand{\xx}{\mathbf{x}}
\renewcommand{\pp}{\mathbf{p}}
\renewcommand{\ggg}{\mathbf{g}}
\newcommand{\GG}{\mathbf{G}}
\DeclareMathOperator{\expc}{\mathbb{E}}
\renewcommand{\Pr}{\mathbb{P}}
\newcommand{\lb}{\underline}
\newcommand{\ub}{\overline}

% mathlcal font
\DeclareFontFamily{U}{dutchcal}{\skewchar\font=45 }
\DeclareFontShape{U}{dutchcal}{m}{n}{<-> s*[1.0] dutchcal-r}{}
\DeclareFontShape{U}{dutchcal}{b}{n}{<-> s*[1.0] dutchcal-b}{}
\DeclareMathAlphabet{\mathlcal}{U}{dutchcal}{m}{n}
\SetMathAlphabet{\mathlcal}{bold}{U}{dutchcal}{b}{n}

% mathscr font (supporting lowercase letters)
%\usepackage[scr=dutchcal]{mathalfa}
%\usepackage[scr=esstix]{mathalfa}
%\usepackage[scr=boondox]{mathalfa}
%\usepackage[scr=boondoxo]{mathalfa}
\usepackage[scr=boondoxupr]{mathalfa}
%\newcommand{\model}{\mathscr{h}}
\newcommand{\model}{\tilde{f}}
\newcommand{\rmod}{F}

\newcommand{\Set}[1]{\mathcal{#1}}
\DeclareMathAlphabet{\mathpzc}{OT1}{pzc}{m}{it} % The mathpzc font
\newcommand{\slv}{\mathpzc}
% mathpzc looks great, but it stops working on 19 Feb 2020 for no reason.
%\newcommand{\slv}{\mathscr}
\newcommand{\software}{\texttt}
\DeclareMathOperator{\eff}{\mathsf{e}\;\!}
\DeclareMathOperator{\Eff}{\mathsf{E}\;\!}
\newcommand{\out}{{\text{out}}}


%% Commands for revision
\newcommand{\red}[1]{\textcolor{red}{#1}}
\newcommand{\blue}[1]{\textcolor{blue}{#1}}
\newcommand{\green}[1]{\textcolor{darkgreen}{#1}}
\newcommand{\TYPO}[1]{{\color{orange}{#1}}}
\newcommand{\MISTAKE}[1]{{\color{violet}{#1}}}
\newcommand{\REPHRASE}[1]{{\color{darkgreen}{#1}}}
\newcommand{\REVISE}[1]{{\color{blue}{#1}}}
\newcommand{\REVISEred}[1]{{\color{red}{#1}}}
\newcommand{\COMMENT}{\todo}  % Needs the todonotes package
%\newcommand{\COMMENT}[1]{\textcolor{brown}{{\small{(comment: #1)}}}}  % This puts comments inline

% Use the following if revision is finished
%\newcommand{\TYPO}{}
%\newcommand{\MISTAKE}{}
%\newcommand{\REPHRASE}{}
%\newcommand{\REVISE}{}
%\newcommand{\REVISEred}{}
%\newcommand{\COMMENT}[1]{}  % Input ignored.


%%%%%%%%%%%%%%%%%%%%%%%%%%%%%%%%%%%%%%%%%%%%%%%%%%%%%%%%%%%%%%%%%%%%%%%%%%%%%%%%%%%%%%%%%%%%%%%%%%%%
\title{ }

\date{\DTMnow}

\author{
    \thanks{}
    %\and
    %Author2
    %\thanks{Information2}
}


\begin{document}

\title{On the Convergence of DFP Algorithm with Wolfe Line Search}
\author{Jinwen Huang}
\date{\today}

\maketitle

%\begin{abstract}
%\end{abstract}

%\textbf{Keywords}: Keyword1, Keyword2
%%%%%%%%%%%%%%%%%%%%%%%%%%%%%%%%%%%%%%%%%%%%%%%%%%%%%%%%%%%%%%%%%%%%%%%%%%%%%%%%%%%%%%%%%%%%%%%%%%%%



%%%%%%%%%%%%%%%%%%%%%%%%%%%%%%%%%%%%%%%%%%%%%%%%%%%%%%%%%%%%%%%%%%%%%%%%%%%%%%%%%%%%%%%%%%%%%%%%%%%%
%% References
% Include references only if there are citations.



\tableofcontents

\section{Introduction}
Global Convergence of the DFP algorithm with Wolfe line search has been a topic of interest in optimization. 
In this article, we will first introduce some background information about the DFP algorithm and the Wolfe line search. 
Then we come to some advancements on the global convergence of the DFP algorithm with Wolfe line search.
Finally, we will list some sufficient conditions for the convergence of the DFP algorithm with Wolfe line search. 



\section{Some background information}

\subsection{DFP Algorithm}
The DFP algorithm is the first quasi-Newton method for unconstrained optimization
\begin{align}
\min_{x \in \RR^n} f(x),
\end{align}
which is first proposed by Davidon\cite{Da} in 1959 and modified by Fletcher and Powell\cite{FP} in 1963.

\begin{algorithm}[H]
\caption{DFP Algorithm}
\begin{algorithmic}[1]
\State Given $x_{1} \in \RR^{n}$; $B_{1} \in \RR^{n \times n}$ positive definite;
    $k:=1$.
\State Compute $g_k = \nabla f(x_k)$;
    If $g_k = 0$ then stop;
    Set $d_k = -B_k^{-1} g_k$.
\State Carry out a search line along $d_k$, getting $\alpha_k > 0$;
    Set $x_{k+1} = x_k + \alpha_k d_k$.
\State Set
    \[B_{k+1} = B_k - \frac{B_{k}s_{k}y_{k}^{T}+y_{k}s_{k}^{T}B_{k}}{s_k^T y_k} - \Bigg(1 + \frac{s_k^T B_k s_k}{s_k^T y_k} \Bigg)\frac{y_k y_k^T}{s_k^T y_k},\]
    where \[s_k = \alpha_{k}d_{k},\]
      \[y_{k} = g_{k+1} - g_{k}.\]
\State $k := k + 1$; go to 2.
\end{algorithmic}
\end{algorithm}

There are two popular methods to obtain $\alpha_k$ in 3. For exact line search, we have $\alpha_k$ satisfies
\[f(x_k + \alpha_k d_k) = \min_{\alpha > 0} f(x_k + \alpha d_k).\]
For inexact line search, we normally requires that 
\[f(x_k + \alpha_k d_k) \leq f(x_k) + c_1 \alpha_k d_k^T g_k,\]
and
\[d_k^T \nabla f(x_k + \alpha_k d_k) \geq c_2 \nabla d_k^T g_k,\]
where $0 < c_1 \leq c_2 < 1$ are constants. Usually $c_1 \leq 0.5$.\\
The first condition is the Armijo condition, which ensures that  $\alpha_k$ leads to a sufficient decrease in the objective function. 
The second condition is the curvature condition, which avoids the step lengths from being too short. 
The Armijo condition and the curvature condition together form the Wolfe conditions.

\begin{remark}
    If we replace the update formula of $B_{k+1}$ in step 4 with
    \[B_{k+1} = B_k - \frac{B_{k}s_{k}s_{k}^{T}B_{k}}{s_k^T B_k s_k} + \frac{y_k y_k^T}{s_k^T y_k},\]
    then we obtain another famous variable metric method, the BFGS algorithm.
\end{remark}

\begin{remark}
The Strong Wolfe conditions are a stronger version of the Wolfe conditions that require the curvature condition to hold with a strict inequality. 
The Strong Wolfe conditions can be expressed as:
\[f(x_k + \alpha d_k) \leq f(x_k) + c_1 \alpha \nabla f(x_k)^T d_k,\]
\[|\nabla f(x_k + \alpha d_k)^T d_k| &\leq c_2 |\nabla f(x_k)^T d_k|,\]
where $c_1$ and $c_2$ are constants in the range $(0, 1)$ with $c_1 < c_2$.
\end{remark}

\subsection{Wolfe Line Search}
In this section, we will analyze the two conditions in Wolfe line search respectively, demonstrating their importance in the iterations.

\subsubsection{Armijo Condition}
The Armijo condition is a sufficient decrease condition that ensures that the step size chosen in the line search leads to a sufficient decrease in the objective function.
The formula for the Armijo condition is given by:
\[f(x_k + \alpha d_k) \leq f(x_k) + c_1 \alpha \nabla f(x_k)^T d_k,\]
where $f$ is the objective function, $x_k$ is the current point, $d_k$ is the search direction, $\alpha$ is the step size, and $c_1$ is a constant in the range $(0, 1)$.  The Armijo condition ensures that the step size chosen in the line search leads to a sufficient decrease in the objective function, which is important for the convergence of optimization algorithms.   

\subsubsection{Curvature Condition}
The curvature condition is a condition that ensures that the step size chosen in the line search is not too small. 
It is given by the formula:
\[\nabla f(x_k + \alpha d_k)^T d_k \geq c_2 \nabla f(x_k)^T d_k,\]
where $c_2$ is a constant in the range $(c_1, 1)$. The curvature condition ensures that the step size chosen in the line search is not too small, which can lead to slow convergence or even divergence of optimization algorithms.



\section{Advancements on the Convergence of the DFP Algorithm}

Powell\cite{Po1} showed that if $f(x)$ is a convex function and if line searches are exact in all iterations, then the DFP algorithm will either stop at a global minimum or generate a sequence of iterates that converges to a global minimum. 
For inexact line search Powell\cite{Po2} showed that the BFGS method is globally convergent. 
Byrd, Nocedal and Yuan\cite{BNY} showed the global convergence of all the methods in the convex Broden family(mentioned below) except the DFP method.\\ 
In the 1990s, Yuan\cite{Yuan}(improved by Yin and Han\cite{YH} three years later), Xu and Liu\cite{XL} gave sufficient conditions for the global convergence of the DFP method with objective function being uniformly convex respectively.
In 2002, Liu, Han and Xu\cite{LHX} proved that convergence of DFP algorithm with Wolfe line search holds for quadratic uniformly convex functions.
In 2024, Luo, Wu and He\cite{LWH} derived a new DFP correction formula and proposed a new class of DFP algorithm under the strong Wolfe condition, improving the problem of linesearch failure due to accuracy and proving global convergence of the improved algorithm under certain assumptions.



\section{Some sufficient conditions for the convergence of the DFP algorithm with Wolfe line search}

\subsection{Yuan's work in the 1990s}




\subsection{Yin ans Han's Improvement}

\subsection{Liu, Han and Xu's work in 2002}

\subsection{Luo, Wu and He's work in 2024}


























\section{Conclusion}
In conclusion, the DFP algorithm and the Wolfe line search are important tools in the field of
optimization. The DFP algorithm is a quasi-Newton method that uses an approximation of the Hessian matrix to find the search direction, while the Wolfe line search is a method for finding an appropriate step size in optimization algorithms. The combination of the DFP algorithm with the Wolfe line search allows for faster convergence and improved performance compared to using the DFP algorithm alone. Recent research has focused on improving the performance and applicability of both the DFP algorithm and the Wolfe line search, making them valuable tools for solving optimization problems in various fields.

\section{References}
%%%
%%%\begin{enumerate}
%%%    \item Y. Yuan. Convergence of DFP algorithm[J]. Science in China, Ser. A, 1995,(11)1281-1294.{p1}
%%%    \item D. XU. and G. Liu., A New Sufficient Condition for the Convergence of DFP Algorithm with Wolfe Linesearch[J].Systems Science and Mathematical Sciences,1996,(03):259-269.{p2}
%%%    \item W. C. Davidon. Variable metric method for minimization, Argonne National Laboratories Report ANL-5990, 1959.  {p3}
%%%    \item R. Fletcher and M. J. D. Powell. A rapidly convergent descent method for minimization, The Computer Journal, 6 (1963),163–168. {p4}
%%%\end{enumerate}

\begin{thebibliography}{99}
    \bibitem[1]{Da}
    W. C. Davidon. 
    Variable metric method for minimization.
    \textit{Argonne National Laboratories Report ANL-5990}, 1959.

    \bibitem[2]{FP}
    R. Fletcher and M. J. D. Powell.
    A rapidly convergent descent method for minimization.
    \textit{The Computer Journal}, 6 (1963), 163–168.

    \bibitem[3]{Po1}
    M. J. D. Powell.
    On the convergence of variable metric algorithm.
    \textit{J. Inst. Math. Appl.}, (1971), 7:21.

    \bibitem[4]{Po2}
    M. J. D. Powell.
    Some Global convergence properties of a variable metric algorithm for minimization without exact line searches.
    \textit{Nonlinear Programming, SIAM-AMS Proceedings vol. IX}, 1976, 53–72.

    \bibitem[5]{BNY}
    R. H. Byrd, J. Nocedal and Y. Yuan.
    Global convergence of a class of variable metric algorithms.
    \textit{SIAM J. Numer. ANal.}, 1987, 24:1171.

    \bibitem[6]{Yuan}
    Y. Yuan.
    On the convergence of the DFP algorithm.
    Technical Report, Computing Center, Academica Sinica, Beijing, China, 1994.

    \bibitem[7]{YH}
    H. Yin and J. Han.
    Some sufficient conditions for the global convergence of the DFP algorithm.
   \textit{Acta Math. Appl. Sin. 21}, No.2, 179-186, 1998.

   \bibitem[8]{XL}
    D. Xu and G. Liu.
    A new sufficient condition for the convergence of DFP algorithm with Wolfe line search.
    \textit{Systems Science and Mathematical Sciences}, 9(1996), 259-269.

    \bibitem[9]{LHX}
    G. Liu, J. Han and D. Xu.
    A note on the convergence of the DFP algoithm on quadratic uniformly convex functions.
    \textit{Optimization 51}, No.2, 339-352, 2002.

    \bibitem[10]{LWH}
    W. Luo, Z. WU and S. He.
    An Improved Quasi-Newtonian Algorithm
    \textit{Journal of Chengdu University of Information Technology}, 2024, 39(03): 374-381.



\end{thebibliography}

\ifnum\value{cite}>0
    \small
    \bibliography{\bibfile}
    \bibliographystyle{plain}
\fi

%% The end
\end{document}
